\begin{lemma}[The two freezing errors, and the moduli of \textup{(4.7)}]
Let $n\ge1$, let $p>1$, let $0<\alpha\le1$, let $C_1\ge0$, let $K>0$, let
$C_\sigma\in\R$, and let $\sigma\colon\R\to\R$ be continuous and nondecreasing
with $\sigma(0)=0$ and $\sigma(t)\le C_\sigma$ for every $t$.  Then there are a
constant $C_2>0$ and functions $\sigma_0,\sigma_1\colon\R\to\R$ which satisfy
\[
\sigma_0(t)=K\bigl(\sigma(t+C_1t^{\alpha})+C_2\,\sigma(t+2C_1t^{\alpha})\bigr),
\qquad
\sigma_1(t)=K\,C_2\,\sigma(t+2C_1t^{\alpha})
\qquad(t>0),
\]
and which are nonnegative, bounded above by one common constant, continuous and
nondecreasing on all of $\R$, with $\sigma_0(0)=\sigma_1(0)=0$, such that the
following holds.

Let $\varepsilon,z_0,R\in\R$ with $R>0$, let $x_0\in\R^n$, let
$f\colon\R^n\times\R\times\R^n\to\R$, let $f^0\colon\R^n\to\R$, let
$u,v\colon\R^n\to\R$ and $g,g_v\colon\R^n\to\R^n$, and write $B=B_R(x_0)$.
Assume
\begin{itemize}
\item $g\in L^p(B)$ and $g_v\in L^p(B)$;
\item each of $x\mapsto f^0(g(x))$, $x\mapsto f^0(g_v(x))$,
  $x\mapsto f(x,u(x),g(x))$ and $x\mapsto f(x,v(x),g_v(x))$ is integrable on
  $B$;
\item for every $x\in B$, every $z\in\R$ and every $h\in\R^n$,
  \[
  \bigl|f(x,z,h)-f^0(h)\bigr|
  \le(1+|h|^2)^{\frac p2}\,\sigma\bigl(|x-x_0|+|z-z_0|\bigr);
  \]
\item $|u(x)-z_0|\le C_1R^{\alpha}$ and $|v(x)-z_0|\le C_1R^{\alpha}$ for
  almost every $x\in B$;
\item $\displaystyle\int_Bf(x,u,g)\,dx\le\int_Bf(x,v,g_v)\,dx$;
\item $\displaystyle\int_B\bigl(\varepsilon+|g|^2+|g_v|^2\bigr)^{\frac p2-1}
  |g-g_v|^2dx\le K\Bigl\{\int_Bf^0(g)\,dx-\int_Bf^0(g_v)\,dx\Bigr\}$.
\end{itemize}
Then
\[
\int_B\bigl(\varepsilon+|g|^2+|g_v|^2\bigr)^{\frac p2-1}|g-g_v|^2dx
\le\sigma_0(R)\int_B(1+|g|^2)^{\frac p2}dx
+\sigma_1(R)\int_B|g-g_v|^{p}\,dx .
\]
\end{lemma}
\begin{proof}
\emph{Step 1: the constants and the two moduli.}  Put
\[
C_2=\max\bigl\{1,2^{\frac p2-1}\bigr\}\cdot2^{\frac p2}>0,
\]
and define, for $t>0$,
\[
\sigma_0(t)=K\bigl(\sigma(t+C_1t^{\alpha})+C_2\,\sigma(t+2C_1t^{\alpha})\bigr),
\qquad
\sigma_1(t)=K\,C_2\,\sigma(t+2C_1t^{\alpha}),
\]
and $\sigma_0(t)=\sigma_1(t)=0$ for $t\le0$.  Note $C_\sigma\ge\sigma(0)=0$, and
that the two displayed formulas for $t>0$ are the two assertions of the
statement.  \emph{These are read off the thesis, not invented}: they are
exactly the two bracketed coefficients the thesis displays in
(4.7) of the thesis at l.~1150--1152, divided by the $\nu_0C$ of (4.5) of the thesis
which the consumer supplies as $K$.  The thesis writes its own explicit
formulas for the same construction one lemma later, at l.~1526--1527 for the
\S5 analogue, and those differ in exactly two legitimate ways: they carry an
extra $2\gamma_4C_1R^{\alpha}$, for which \S4 has no $\gamma_4$, and they use
$C_1$ where these use $2C_1$, because \S5 compares against a fixed point of the
sphere rather than against the centre value.

For $c\ge0$ write $S_c(t)=t+c\,t^{\alpha}$.  Since $\alpha>0$, the real power
$t\mapsto t^{\alpha}$ is continuous at every real $t$, so $S_c$ is continuous on
$\R$; $S_c(0)=0+c\cdot0^{\alpha}=0$ because $0^{\alpha}=0$ for $\alpha\ne0$;
$S_c(t)\ge0$ for $t\ge0$, since both summands are; and $S_c$ is nondecreasing on
$[0,\infty)$, because $t\mapsto t^{\alpha}$ is nondecreasing there.

\emph{Nonnegativity.}  For $t\le0$ both functions are $0$.  For $t>0$ we have
$S_{C_1}(t)\ge0$ and $S_{2C_1}(t)\ge0$, so
$\sigma(S_{C_1}(t))\ge\sigma(0)=0$ and $\sigma(S_{2C_1}(t))\ge0$ by
monotonicity, and $K,C_2>0$.

\emph{Boundedness.}  Take $C=K(1+C_2)C_\sigma$.  For $t\le0$ both functions are
$0\le C$, because $K,C_2>0$ and $C_\sigma\ge0$.  For $t>0$,
$\sigma_0(t)\le K(C_\sigma+C_2C_\sigma)=C$ and
$\sigma_1(t)\le KC_2C_\sigma\le C$.

\emph{Continuity.}  On $(-\infty,0]$ each function is the constant $0$, and on
$(0,\infty)$ each is a composition and a linear combination of continuous
functions, $\sigma$ and $S_c$ both being continuous.  At $t=0$ the two
prescriptions agree: $S_{C_1}(0)=S_{2C_1}(0)=0$ and $\sigma(0)=0$, so the
second prescription also returns $0$ there.  Hence both functions are
continuous on $\R$.

\emph{Monotonicity.}  Let $s\le t$.  If $t\le0$ both values are $0$.  If
$s\le0<t$ the value at $s$ is $0$ and the value at $t$ is nonnegative, as just
shown.  If $0<s\le t$ then $S_{C_1}(s)\le S_{C_1}(t)$ and
$S_{2C_1}(s)\le S_{2C_1}(t)$, so the same inequalities hold after applying the
nondecreasing $\sigma$, and then after multiplying by the positive $K$ and
$C_2$.

\emph{Value at $0$.}  Both are $0$ by definition, $0\le0$.

Now fix the data, with $R>0$; since $R>0$ we have $\sigma_0(R)=
K(\sigma(R+C_1R^{\alpha})+C_2\sigma(R+2C_1R^{\alpha}))$ and
$\sigma_1(R)=KC_2\sigma(R+2C_1R^{\alpha})$.

\emph{Step 2: two integrability facts.}  First, for $s\ge0$,
$1+s^2\le(1+s)^2$, so $(1+s^2)^{p/2}\le(1+s)^{p}$ because $t\mapsto t^{p/2}$ is
nondecreasing on $[0,\infty)$ and $((1+s)^2)^{p/2}=(1+s)^p$.  If
$h\in L^p(B)$ is $\R^n$-valued then $1+|h|\in L^p(B)$, the constant $1$ lying
in $L^p(B)$ because $|B|<\infty$ (Preamble: the volume of a ball), so
$(1+|h|)^p\in L^1(B)$; and $x\mapsto(1+|h(x)|^2)^{p/2}$ is measurable, being a
continuous function of $|h|$.  By domination (Preamble),
$x\mapsto(1+|h(x)|^2)^{p/2}$ is integrable on $B$.  We use this for $h=g$ and
$h=g_v$.  Second, $g-g_v\in L^p(B)$, so $|g-g_v|^p$ is integrable on $B$.

\emph{Step 3: the elementary power inequality and \textup{(6.5)}.}  For
$q>0$ and $X,Y\ge0$,
\[
(X+Y)^q\le\max\{1,2^{q-1}\}\,(X^q+Y^q),
\tag{3a}
\]
by subadditivity of $t\mapsto t^q$ for $0<q\le1$ and by the power-mean
inequality for $q\ge1$.  Now fix $x\in B$.  From $g_v=g-(g-g_v)$ we get
$|g_v|\le|g|+|g-g_v|$, hence
$|g_v|^2\le2|g|^2+2|g-g_v|^2$ and so
\[
1+|g_v(x)|^2\le2\bigl(1+|g(x)|^2\bigr)+2|g(x)-g_v(x)|^2 .
\]
Raising to the power $p/2\ge0$, which is nondecreasing on $[0,\infty)$, and
applying (3a) with $q=p/2$, $X=2(1+|g|^2)$ and $Y=2|g-g_v|^2$, and then
$(2Z)^{p/2}=2^{p/2}Z^{p/2}$ and $(|g-g_v|^2)^{p/2}=|g-g_v|^p$,
\[
\bigl(1+|g_v(x)|^2\bigr)^{\frac p2}
\le C_2\Bigl(\bigl(1+|g(x)|^2\bigr)^{\frac p2}+|g(x)-g_v(x)|^{p}\Bigr).
\]
Both sides are integrable on $B$ by Step 2, so integrating,
\[
\int_B(1+|g_v|^2)^{\frac p2}
\le C_2\Bigl\{\int_B(1+|g|^2)^{\frac p2}+\int_B|g-g_v|^{p}\Bigr\}.
\tag{3b}
\]

\emph{Step 4: each freezing error.}  Let $w\colon\R^n\to\R$ and
$h\colon\R^n\to\R^n$, let $c\ge0$, suppose $x\mapsto f(x,w(x),h(x))$,
$x\mapsto f^0(h(x))$ and $x\mapsto(1+|h|^2)^{p/2}$ are integrable on $B$ and
that $|w-z_0|\le c$ almost everywhere on $B$.  Then
\[
\Bigl|\int_Bf(x,w,h)\,dx-\int_Bf^0(h)\,dx\Bigr|
\le\sigma(R+c)\int_B(1+|h|^2)^{\frac p2}dx .
\tag{4a}
\]
Indeed, the difference of the two integrals is the integral of the difference,
which is integrable; the absolute value of an integral is at most the integral
of the absolute value (Preamble); and for almost every $x\in B$ the third
hypothesis of the lemma at $(x,w(x),h(x))$ gives
\[
\bigl|f(x,w(x),h(x))-f^0(h(x))\bigr|
\le(1+|h(x)|^2)^{\frac p2}\sigma\bigl(|x-x_0|+|w(x)-z_0|\bigr)
\le(1+|h(x)|^2)^{\frac p2}\sigma(R+c),
\]
the second step because $|x-x_0|\le R$ on $B$, $|w(x)-z_0|\le c$ almost
everywhere, and $\sigma$ is nondecreasing, the weight being nonnegative.
Monotonicity of the integral (Preamble) gives (4a).

Apply (4a) twice: with $(w,h,c)=(u,g,C_1R^{\alpha})$, and with
$(w,h,c)=(v,g_v,2C_1R^{\alpha})$ --- legitimate because
$|v-z_0|\le C_1R^{\alpha}\le2C_1R^{\alpha}$ almost everywhere, $C_1R^{\alpha}$
being nonnegative.  This gives
\[
\int_Bf^0(g)-\int_Bf(x,u,g)\le\sigma(R+C_1R^{\alpha})\int_B(1+|g|^2)^{\frac p2},
\qquad
\int_Bf(x,v,g_v)-\int_Bf^0(g_v)
\le\sigma(R+2C_1R^{\alpha})\int_B(1+|g_v|^2)^{\frac p2},
\]
each being at most the corresponding absolute value.

\emph{Step 5: \textup{(4.7)} and the conclusion.}  Writing the four integrals
as real numbers and adding the two displays of Step 4 to the minimality
hypothesis $\int_Bf(x,u,g)\le\int_Bf(x,v,g_v)$,
\[
\int_Bf^0(g)-\int_Bf^0(g_v)
\le\sigma(R+C_1R^{\alpha})\int_B(1+|g|^2)^{\frac p2}
+\sigma(R+2C_1R^{\alpha})\int_B(1+|g_v|^2)^{\frac p2},
\]
because the left-hand side is the sum of the three brackets
$\bigl\{\int_Bf^0(g)-\int_Bf(x,u,g)\bigr\}$,
$\bigl\{\int_Bf(x,v,g_v)-\int_Bf^0(g_v)\bigr\}$ and
$\bigl\{\int_Bf(x,u,g)-\int_Bf(x,v,g_v)\bigr\}$, of which the third is at most
$0$.  Substituting (3b) into the second term --- legitimate because
$\sigma(R+2C_1R^{\alpha})\ge\sigma(0)=0$, the argument being nonnegative ---
\[
\int_Bf^0(g)-\int_Bf^0(g_v)
\le\bigl\{\sigma(R+C_1R^{\alpha})+C_2\sigma(R+2C_1R^{\alpha})\bigr\}
\int_B(1+|g|^2)^{\frac p2}
+C_2\sigma(R+2C_1R^{\alpha})\int_B|g-g_v|^{p} ,
\]
which is the thesis' (4.7) of the thesis.  Multiplying by $K>0$ and combining with
the last hypothesis of the lemma,
\[
\int_B\bigl(\varepsilon+|g|^2+|g_v|^2\bigr)^{\frac p2-1}|g-g_v|^2
\le K\Bigl\{\int_Bf^0(g)-\int_Bf^0(g_v)\Bigr\}
\le\sigma_0(R)\int_B(1+|g|^2)^{\frac p2}+\sigma_1(R)\int_B|g-g_v|^{p},
\]
by the formulas for $\sigma_0(R)$ and $\sigma_1(R)$ recorded at the end of
Step 1.  That is the assertion.
\end{proof}
